\documentclass{article}
\usepackage[UTF8]{ctex}
\usepackage{amsmath}
\usepackage{amssymb}

\begin{document}

\section*{复域方程 ((x^5 + 2)/x)^((x^5 + 3)/x) = x^(x - 1) 的严谨解析}

\subsection*{1. 奇点与分枝切割}
方程在 $x=0$ 处有本质奇点, 在 $x^5+2=0$ (即 $x = \sqrt[5]{2} e^{i\frac{(2n+1)\pi}{5}}$) 处有分歧点. 
令 $x = r e^{i\theta}$, 代入复对数展开式:
\[ (x^5+3) \left[ \ln|x^5+2| + i\phi_k \right] = (x^5+x^2-x+3) \left[ \ln r + i(\theta + 2m\pi) \right] \]
其中 $\phi_k = \text{Arg}(x^5+2) + 2k\pi$. 

\subsection*{2. 虚数解的离散分布 (Spectral Distribution)}
考察当 $|x| \to 1$ 时的渐近行为. 令 $x = e^{i\theta}$. 
此时 $|x^5+2| = |e^{i5\theta} + 2| = \sqrt{5 + 4\cos(5\theta)}$. 
代入虚部等式 (Imaginary Part Equality):
\[ \text{Im} \left[ (e^{i5\theta}+3) \ln|e^{i5\theta}+2| \right] = \text{Im} \left[ (e^{i5\theta}+e^{i2\theta}-e^{i\theta}+3) i(\theta + 2m\pi) \right] \]
利用欧拉公式展开, 当 $\theta = \pi$ 时:
\[ x = e^{i\pi} = -1 \]
左侧: \((-1+3) \ln|-1+2| = 2 \ln 1 = 0\). 
右侧: \((-1+1-(-1)+3) i(\pi + 2m\pi) = 4 i\pi(1+2m)\). 
由于原方程是指数形式 $A^B = e^{B \text{Log} A}$, 
左侧为 $e^0 = 1$, 右侧为 $e^{4i\pi(1+2m)} = (e^{2i\pi})^{2(1+2m)} = 1$. 
这严谨地证明了 $x = -1$ 是所有分枝下的共同基点解.

\subsection*{3. 超越虚根的 Lambert 映射}
对于 $|x| \neq 1$ 的纯虚数解 $x = iy$, 方程退化为:
\[ (i y^5 + 3) \ln(i y^5 + 2) = (i y^5 - y^2 - iy + 3) (\ln |y| + i\frac{\pi}{2}) \]
通过实部与虚部分离, 得到关于 $y$ 的超越方程组. 该解集对应于黎曼面上的一组离散点, 其渐近分布满足:
\[ x_n \approx \pm i \exp \left( \frac{2n\pi}{ \text{Re}(\text{Residue}) } \right) \]

\end{document}